\begin{abstract}
Corporate leverage varies enormously across firms, and existing theories account for only a fraction of the observed dispersion. This paper argues that a key missing ingredient is the \emph{shape} of the capital supply distribution: how lender endowments align with a firm's cash flow distribution determines borrowing capacity in ways that mean and variance alone cannot capture. I formulate debt contracting as an optimal transport problem and show that the Wasserstein distance between a firm's cash flow distribution and the aggregate capital supply distribution governs both leverage and credit spreads. Within-firm panel regressions confirm that greater distributional mismatch reduces leverage, controlling for standard determinants; the effect is strongest among bond-issuing firms, where the capital supply distribution governs pricing most directly. The model delivers a new welfare measure for financial intermediation: aggregate deadweight bankruptcy costs equal the squared two-Wasserstein distance between the distribution of firm cash flows and the capital supply distribution.
\end{abstract}
